\documentclass[oneside]{scrbook}

% \usepackage[utf8]{inputenc}
\usepackage[T1]{fontenc}
\usepackage{version}
% \usepackage[latin1]{inputenc}
\usepackage[italian]{babel}
\usepackage{algpseudocode}
\usepackage{amsmath}
\usepackage{booktabs}
\usepackage{color}
\usepackage{amsfonts}
\usepackage{url}
\usepackage{algorithm}
\usepackage{chngcntr}
\usepackage[answerdelayed]{exercise}
\usepackage{graphicx}
\usepackage{caption}
\usepackage{listings}
\usepackage{tikz}
\usepackage{pgf}
\usepackage{colortbl}
\usepackage{listingsutf8}
\usepackage[pdftitle = {Dispense di esercizi}, pdfauthor = {Gianmarco Tedeschi e Alberto Marchesi}, pdfsubject = {Fondamenti di Informatica ($082746$ - A.A.~$2025-2026$)}, pdfcreator = {Gianmarco Tedeschi e Alberto Marchesi}, pdfproducer = {-}]{hyperref}
% \usepackage[automark,nouppercase]{scrpage2}
\usepackage{mdframed}
\usepackage{scrlayer-scrpage} % Necessary package for page styles

%\pagestyle{scrheadings}

%\usepackage{titlesec}
%\newcommand{\sectionbreak}{\clearpage}
\counterwithin{Exercise}{chapter}
\counterwithin{Answer}{chapter}

\algrenewcommand\algorithmicthen{\textbf{Allora}}
\algrenewcommand\algorithmicif{\textbf{Se}}
\algrenewcommand\algorithmicelse{\textbf{Altrimenti}}
\algrenewcommand\algorithmicend{\textbf{Chiudi}}
\algrenewcommand\algorithmicwhile{\textbf{Finch\'e}}
\algrenewcommand\algorithmicdo{\textbf{Esegui}}
\algrenewcommand{\algorithmiccomment}[1]{\hskip3em\texttt{//#1}}

\setlength{\parindent}{0in}
\setlength{\parskip}{10pt}

\newmdenv[%
        outerlinewidth=2,%
        roundcorner=10pt,%
        leftmargin=0,%
        rightmargin=0,%
        outerlinecolor=black,
        innertopmargin=\topskip,%
        splittopskip=\topskip,%
]{algo}


\ohead{\headmark}
\ihead{}
\chead{}
\ifoot{A.A.~$2025-2026$}
\cfoot{Fondamenti di Informatica (082746) - G. Tedeschi}
\ofoot{Pag.~\pagemark}
\setheadsepline{.4pt}
\setfootsepline{.4pt}

\definecolor{lbcolor}{rgb}{0.9,0.9,0.9}

\lstdefinestyle{cstyle}{%
  inputencoding=utf8/latin1,
  breaklines=true,
  breakatwhitespace=false,
  basicstyle=\small\ttfamily,
  language=c,
  frame=single,
  showstringspaces=false
}

\lstdefinestyle{cstylesmall}{%
  inputencoding=utf8/latin1,
  breaklines=true,
  breakatwhitespace=false,
  basicstyle=\scriptsize\ttfamily,
  language=c,
  frame=single,
  showstringspaces=false
}

\lstdefinestyle{console}{%
  inputencoding=utf8/latin1,
  breaklines=true,
  breakatwhitespace=false,
  basicstyle=\scriptsize\ttfamily,
  showstringspaces=false,
  backgroundcolor=\color{lbcolor}
}


\lstnewenvironment{clang}{\lstset{style=cstyle}}{}
\lstnewenvironment{clangsmall}{\lstset{style=cstylesmall}}{}

\addtocounter{chapter}{3}

\begin{document}
\chapter{Array numerici}
Questa dispensa propone esercizi sulla scrittura di algoritmi, in linguaggio C, utili alla comprensione dei vettori (array). 

La dichiarazione di un vettore di elementi omogenei in C avviene grazie alla seguente dichiarazione:
\begin{lstlisting}[style=cstyle]
void main(){

tipo_variabile nome_variabile[numero_elementi];

istruzioni;

}
\end{lstlisting}
dove \texttt{tipo\_variabile} \`e il tipo degli elementi del vettore e \texttt{numero\_elementi}  \`e la lunghezza del vettore (considerata una costante da qui in poi). Solitamente la lunghezza del vettore viene dichiarato come una costante tramite la parola chiave \texttt{\#define} per migliorare la leggibilit\`a del codice.

Per accedere a singoli elementi dell'array utilizzeremo degli indici che varieranno tra $0$ e \texttt{numero\_elementi} - $1$. Per accedere ad un elemento del vettore si utilizzano le parentesi quadre, ad esempio:
\begin{lstlisting}[style=cstyle]
void main(){

int vettore[10];

vettore[1] = 4;

}
\end{lstlisting}
il precedente programma dichiara un vettore di interi di dimensione $10$ e successivamente inserisce il valore $4$ nel secondo elemento del vettore.

Un vettore ha dimensioni fisse. Esse devono essere specificate alla dichiarazione del programma. Non \`e possibile modificare le dimensioni di un vettore durante l'esecuzione del programma.

\clearpage
\section{Esercizi}

% Facili
\begin{Exercise}[label={ex:inversione}]
  Scrivere un programma che permetta all'utente di inserire un array di numeri v1 (al più 100 valori). Il programma copia l'array v1 al contrario nell'array v2. Il programma stampa v2.
\end{Exercise}

\begin{Answer}[ref={ex:inversione}]
  \lstinputlisting[style=cstylesmall]{src/inversione.c}
  \paragraph{Nota: } Alternativamente, possiamo copiare e stampare il vettore invertito in un unico ciclo.
  \lstinputlisting[style=cstylesmall]{src/inversione_2.c}
\end{Answer}

\begin{Exercise}[label={ex:coppie}]
  Scrivere un programma che permetta all'utente di inserire una sequenza di numeri. Il programma cerca successivamente le coppie di numeri per cui il primo è il doppio del secondo e le stampa.
\end{Exercise}

\begin{Answer}[ref={ex:coppie}]
  \lstinputlisting[style=cstylesmall]{src/coppie.c}
\end{Answer}

\begin{Exercise}[label={ex:pari_dispari}]
		Scrivere un programma che acquisisca un certo numero di valori interi e positivi stabilito dall'utente (massimo $10$), e poi stampi a video un vettore contenente tutti i numeri pari della sequenza e un vettore contente tutti i numeri dispari.
\end{Exercise}

\begin{Answer}[ref={ex:pari_dispari}]
  \lstinputlisting[style=cstylesmall]{src/pari_dispari.c}
\end{Answer}
\begin{Exercise}[label={ex:uguaglianza}]
  Scrivere un programma che legga da input un numero non noto a priori di interi positivi (al più 10) terminati da 0. Qualora l'utente inserisca dei valori negativi, tali valori dovranno essere scartati e non considerati. Si memorizzi tale sequenza in un vettore di opportuna dimensione. Si stampino a video tutti i numeri che sono allocati nel vettore in posizioni il cui indice è uguale al numero stesso.
\end{Exercise}

\begin{Answer}[ref={ex:uguaglianza}]
  \lstinputlisting[style=cstylesmall]{src/uguaglianza.c}
\end{Answer}

\begin{Exercise}[label={ex:indici}]
  Scrivere un programma che acquisisca un certo numero di valori stabilito dall'utente (massimo $10$), e poi acquisisca un eguale quantit\`a di indici, che specificano l'ordine in cui il programma dovr\`a stampare i valori acquisiti in precedenza. Per esempio:

\begin{lstlisting}[style=console]
Quanti elementi vuoi inserire (max: 10)? 5
Inserire il 1o valore: 10
Inserire il 2o valore: 20
Inserire il 3o valore: 30
Inserire il 4o valore: 40
Inserire il 5o valore: 50
Qual è il 1o valore che vuoi stampare? 3
Qual è il 2o valore che vuoi stampare? 2
Qual è il 3o valore che vuoi stampare? 1
Qual è il 4o valore che vuoi stampare? 4
Qual è il 5o valore che vuoi stampare? 5
30
20
10
40
50
\end{lstlisting}
\end{Exercise}

\begin{Answer}[ref={ex:indici}]
  \lstinputlisting[style=cstylesmall]{src/indice.c}
\end{Answer}

\begin{Exercise}[label={ex:massimo}]
  Scrivere un programma che chieda all'utente di inserire un numero variabile (al massimo $10$) di valori reali. Il programma dovr\`a assicurarsi che l'utente non possa inserire pi\`u di $10$ valori.

  Il programma dovr\`a cercare, durante l'acquisizione, la posizione del minimo e del massimo.

  Infine, il programma dovr\`a stampare un sommario dei valori inseriti, con l'indicazione della posizione del minimo (indicata con un segno $-$) e del massimo (indicata con un segno $+$), il tutto formattato come indicato di seguito.

\begin{lstlisting}[style=console]
Quanti elementi vuoi inserire (max: 10)? 100
Il valore inserito è troppo grande!
Quanti elementi vuoi inserire (max: 10)? 4
Inserire il 1o valore: 100
Inserire il 2o valore: 20
Inserire il 3o valore: -1
Inserire il 4o valore: 5.3
  +| 100.00
   | 20.00
-  | -1.00
   | 5.30
\end{lstlisting}
\end{Exercise}

\begin{Answer}[ref={ex:massimo}]
  \lstinputlisting[style=cstylesmall]{src/massimo.c}
\end{Answer}


\begin{Exercise}[label={ex:set}]
  Scrivere un programma che acquisisce un numero variabile di interi, con un massimo di $100$. Il programma dovr\`a costruire un insieme (ossia controllare che non ci siano elementi ripetuti) a partire dagli elementi inseriti.
\end{Exercise}

\begin{Answer}[ref={ex:set}]
  \lstinputlisting[style=cstylesmall]{src/set.c}
\end{Answer}

\begin{Exercise}[label={ex:operazioni_set}]
  Implementare le operazioni di intersezione, unione, e differenza insiemistica tra due array acquisiti da tastiera, supponendo siano insiemi (ossia non abbiano elementi ripetuti).
\end{Exercise}

\begin{Answer}[ref={ex:operazioni_set}]
  \lstinputlisting[style=cstylesmall]{src/operazioni-set.c}
\end{Answer}


\begin{Exercise}[label={ex:mischia}]
	Scrivere un programma che, dati due vettori di interi di dimensione $10$, ne costruisca un terzo di dimensione $20$ i cui elementi di posizione pari siano gli elementi del primo vettore e gli elementi di posizione dispari siano gli elementi del secondo vettore.

	Provare ad utilizzare un solo vettore di dimensione $20$ per risolvere l'esercizio.
\end{Exercise}

\begin{Answer}[ref={ex:mischia}]
  \lstinputlisting[style=cstylesmall]{src/mischia.c}
\end{Answer}


\begin{Exercise}[label={ex:media_varianza}]
  	Scrivere un programma che acquisisca un numero stabilito dall'utente di voti di esami e si assicuri che tutti i voti inseriti siano nell'intervallo $[18,30]$. Il numero massimo di voti inseribili \`e $20$, ma l'utente potrebbe decidere di inserire meno voti. Il programma dovr\`a poi stampare a video un sommario simile a quello di seguito. La media troncata \`e la media calcolata non contando gli estremi (voti minimo e massimo).

\begin{lstlisting}[style=console]
STATISTICHE VOTI:

Numero esami sostenuti: ...
Media: ...
Media troncata: ...
Varianza: ...
Deviazione standard: ...
\end{lstlisting}

Per il calcolo la varianza, $\hat{\sigma}^2$, utilizzare la seguente formula:
\begin{equation}
  \hat{\sigma}^2 = \frac{\sum_{i = 1}^{N} (v_i - \hat{v})^2}{N-1}
\end{equation}
dove $v_i$ \`e l'$i$-esimo voto e $\hat{v}$ \`e la media campionaria dei voti. La deviazione standard \`e semplicemente la radice quadrata della varianza.
\end{Exercise}

\begin{Answer}[ref={ex:media_varianza}]
  \lstinputlisting[style=cstylesmall]{src/media_varianza.c}
\end{Answer}

\begin{Exercise}[label={ex:isto_divisori}]
	Scrivere un programma che, acquisita una sequenza di al massimo $50$ numeri interi tra $1$ e $100$, stampi a video l'istogramma dei divisori. L'istogramma deve avere tutti i valori da $2$ al valore massimo immesso diviso per $2$. Il programma deve considerare solo i divisori propri, dove un divisore positivo di $n$ diverso da $n$ stesso e  $1$ \`e chiamato divisore proprio.

	Un esempio dell'output del programma \`e:
\begin{lstlisting}[style=console]
Quanti numeri vuoi considerare: 5

Inserire il 1o numero: 6
Inserire il 2o numero: 4
Inserire il 3o numero: 3
Inserire il 4o numero: 9
Inserire il 5o numero: 12
2		| ***
3		| ***
4		| *
5		| 
6		| *
\end{lstlisting}
\end{Exercise}

\begin{Answer}[ref={ex:isto_divisori}]
  \lstinputlisting[style=cstylesmall]{src/isto_divisori.c}
\end{Answer}



\clearpage
\section*{Soluzioni}
\shipoutAnswer

\end{document}